\chapter{Dynamics}
\section{First law}
\begin{definition}{N1L}{n1l}
	A body will continue in its state of rest or move at constant speed in a straight line,
	unless an \it{external resultant force} acts on it.
\end{definition}
\begin{definition}{Inertia}{inertia}
	The resistance to a change in the state of motion of an object.
\end{definition}
\section{Second law}
\begin{definition}{N2L}{n2l}
	The \it{rate of change of linear momentum} of a body is \it{directly proportional to the resultant force} acting on it
	and its direction is in the \it{same direction} as this resultant force.
	\[\sum \va{F} \propto \dot{\va{p}} = \odv{\ab(mv)}{t}\]
\end{definition}
When a body has constant mass, \cref{def:n2l} can be simplified to
\begin{equation*}
	\sum \va{F} = m\dot{v} = ma
	\label{eq:n2l}
\end{equation*}
\section{Third law}
\begin{definition}{N3L}{n3l}
	If body A exerts a force on body B, then body B exerts an \it{equal and opposite} force on body A.
	\[\va{F}_\text{AB} = -\va{F}_\text{BA}\]
	This results in an action-reaction pair, for which
	\begin{itemize}
		\item forces act on different bodies,
		\item are of the same nature,
		\item and are equal in magnitude and opposite in direction
	\end{itemize}
\end{definition}
\section{Momentum}
\begin{definition}{PCOLM}{pcolm}
	The \it{total momentum} of a closed system is constant when no external resultant force acts on it.
	\[\sum p_i = \sum p_f\]
\end{definition}

\begin{definition}{Elasticity}{elasticity}
	A collision between two objects is \bf{perfectly elastic} if the total kinetic energy of the system before and after
	the collision is unchanged.
\end{definition}
For a perfectly elastic collision between objects A and B, \(\va{v}_A - \va{v}_B = \va{u}_B - \va{u}_A\); the relative
speed of approach equals the relative speed of separation.

\begin{example}{}{pcolm}
	A tritium nucleus of mass \(3m\) and initial speed \(u\) collides \it{elastically} with a deuterium nucleus
	of mass \(2m\) and initial speed \(-u\). We can solve for the final velocities \(v_D\) and \(v_T\) by applying
	\cref{def:pcolm,def:elasticity}:
	\begin{align*}
		3mu - 2mu & = 3mv_T + 2mv_D \\
		u - (-u)  & = v_D - v_T
	\end{align*}
\end{example}

\begin{definition}{Impulse}{impulse}
	The product of the \it{average force} acting on an object and the \it{time interval} during which it acts.
	\[\va{J} = \ab<F>\Delta t = \int \va{F}\odif{t}\]

	The \bf{impulse-momentum theorem} states that the impulse of a force acting on an object is equal to its change in momentum.
	\[\va{J} = \Delta \va{p} = \va{p}_f-\va{p}_i\]
\end{definition}

\section{Special types of forces}
\subsection{Elastic spring force}
For elastic strings or springs that obey \bf{Hooke's Law}, the
external force pulling on the string is
\begin{equation*}
	F = k\Delta x = k\ab(l - l_0)
\end{equation*}
where \(k\) is the \bf{spring constant} and \(x\) is the \bf{extension} of the spring.
The elastic potential energy is the area under the \(F\)-\(x\) graph:
\begin{equation*}
	U = \f{1}{2}kx^2
\end{equation*}
\subsection{Upthrust}
\begin{definition}{Upthrust}{upthrust}
	The \it{net upward force} exerted by a fluid on a body floating/submerged in it.
	\[U = m_f g = \rho_f V_f g\]
	where \(m_f\) is the mass of fluid displaced, \(\rho_f\) is the density of the fluid and \(V_f\) is the volume of fluid displaced.
\end{definition}
\subsection{Drag force}
\begin{definition}{Drag force}{drag-force}
	The force resisting an object \it{moving relative to a fluid} that \it{opposes}
	its motion.
\end{definition}
